% \begin{abstract}
% Vision transformers retain fine image detail but pay to process every patch
% through every layer. Token merging reduces this cost only when a trainable
% merging policy also produces the shorter tensor used by downstream blocks.
% We introduce \mn{}, a local-to-latent vision transformer that separates these
% roles: continuous spatial mass routing learns where information should flow,
% and an exact-budget gather physically reduces $N$ patches to $K$ carriers
% before global reasoning. Routing is restricted to the original image grid,
% but similarities, transported mass, and carrier locations remain learned. A
% threshold-based soft top-$k$ relaxation trains the routing scores, while a
% straight-through carrier weight connects them to the hard gather. Selected
% mass is retained as a logarithmic key bias; we also give an exact rank-one
% query/key augmentation that implements this bias with an unmodified
% FlashAttention kernel. On ImageNet-1K with DeiT-S/8, the trained checkpoint
% reaches 81.368\% top-1 at $2\times$ carrier compression, versus 82.246\% for
% the dense model; matched-budget evaluation gives 81.360\% for \mn{} and
% 81.932\% for post-training ToMe. Radius-3 routing improves a matched
% 150-epoch control by 1.670 percentage points over global routing and 0.598
% over flattened windows, and three-seed CIFAR-100 controls retain the spatial
% advantage after matching receiver degree. The remaining
% \pendingresult{DTEM score}, \pendingresult{latency}, and
% \pendingresult{memory} cells are explicitly reserved for one common-protocol
% checkpoint evaluation.
% \end{abstract}

\begin{abstract}
Modern visual recognition requires models to capture both fine-grained image
detail and long-range context. Vision transformers meet these needs by
representing images as patch tokens and modeling their interactions, but
retaining every token across layers becomes costly at high resolution. Token
merging reduces this cost by consolidating redundant tokens. Yet its success
depends both on learning which tokens to combine without losing task-relevant
information and on translating those decisions into a physically shorter
sequence. DTEM learns the merging policy with dedicated merge embeddings and a
continuous relaxation. However, its soft training path retains every token
slot: it learns merge decisions without realizing physical compression, which
still requires hard merging. \mn{} closes this gap with a local-to-latent
architecture. Differentiable routing consolidates information according to the
learned merge decisions; fixed-budget selection then physically shortens the
sequence before global reasoning. To make this transition trainable, we combine
differentiable budgeted selection with mass-aware attention over merged tokens.
We evaluate \mn{} on ImageNet-1K image classification, using DeiT-S/8 as the
dense baseline and comparing with ToMe and PiToMe at matched token budgets.
At a $2\times$ carrier reduction, the trained
\mn{} checkpoint obtains 81.360\% top-1, compared with 82.246\% for dense
DeiT-S/8, 81.932\% for post-training ToMe, and 81.742\% for PiToMe; the matched
DTEM-p8 adaptation reaches \pendingvalue\%. Among the 150-epoch ImageNet routing
variants, radius-3 routing scores 1.670 percentage points above global routing
and 0.598 points above flattened-window routing. A three-seed CIFAR-100 study
also favors radius-3 routing over global, flattened-window, and degree-matched
receiver controls. These results characterize the spatial routing variant
and physical bottleneck. Under the common trained-checkpoint protocol, the
accuracy--efficiency comparison is \pendingresult{comparison summary}.
\end{abstract}
